
\section{Methods}\label{sec:general_method}
\subsection{Human Experiment Data Collection}

\paragraph{Reconstructing human bid distributions via moment matching.}
Because micro-level bid data are rarely available from published experiments, we reconstruct synthetic human bid distributions by calibrating a mixture model to aggregate statistics
reported in each source. We model bids as:
\[
b = b^\star_m(v) \cdot r
\]
where $b^\star_m(v)$ is the equilibrium bid and $r$ is a bid-to-equilibrium ratio drawn from a three-component mixture:
\begin{align*}
\text{With probability } p_{\text{eq}}: \quad & r \sim \mathcal{N}(1, \sigma^2) \\
\text{With probability } p_{\text{over}}: \quad & r \sim 1 + \text{Exp}(\lambda) \\
\text{With probability } p_{\text{under}}: \quad & r \sim \text{Uniform}(1-\delta, 1)
\end{align*}
This structure maps directly to how experimental papers categorize behavior: equilibrium bidders ($r \approx 1$), overbidders ($r > 1$), and underbidders ($r < 1$).

\paragraph{Source papers and equilibrium benchmarks.}
We draw on four sources: \citet{li2017obviously} and \citet{breitmoser2022obviousness} for second-price and ascending clock auctions; \citet{kagel1993independent} for first-price,
second-price, and third-price sealed-bid auctions with $n \in \{5, 10\}$ bidders; and \citet{gonczarowski2023strategyproofness} for second-price auctions comparing mechanism
descriptions. For each format, we use the risk-neutral Bayes-Nash equilibrium:
\begin{align*}
b^\star_{\text{SPSB}}(v) &= v & &\text{(second-price, ascending clock)} \\
b^\star_{\text{FPSB}}(v) &= \tfrac{n-1}{n} v & &\text{(first-price)} \\
b^\star_{\text{TPSB}}(v) &= \tfrac{n-1}{n-2} v & &\text{(third-price)}
\end{align*}

\paragraph{Parameter identification.}
Mixture weights $(p_{\text{eq}}, p_{\text{over}}, p_{\text{under}})$ are set directly from reported rates. For example, \citet{li2017obviously} reports ${\sim}50\%$ dominant-strategy
play and ${\sim}40\%$ overbidding, yielding $p_{\text{eq}} = 0.50$, $p_{\text{over}} = 0.40$, $p_{\text{under}} = 0.10$.

The overbid scale $\lambda$ is calibrated to match a secondary moment:
\begin{itemize}
  \item For \citet{li2017obviously} and \citet{breitmoser2022obviousness}: the mean bid-to-value ratio (e.g., $\mathbb{E}[b/v] \approx 1.08$).
  \item For \citet{kagel1993independent}: the regression coefficient $\beta$ from bid regressions.
  \item For \citet{gonczarowski2023strategyproofness}: mean absolute deviation (MAD $\approx$ \$0.51--0.55).
\end{itemize}
The equilibrium noise $\sigma$ is derived from regression $R^2$ where reported, via $\sigma^2 = \beta^2 \text{Var}(v)(1-R^2)/R^2$. The underbid bound $\delta$ is set consistently across
papers; this is a parsimony choice as papers do not report underbid tail shapes separately.

\paragraph{Why mixture model over QRE.}
Quantal Response Equilibrium predicts symmetric deviations around equilibrium, contradicting the well-documented asymmetric overbidding in second-price auctions
\citep{kagel1993independent}. The mixture model directly maps to how experimental papers categorize and report behavior.

\paragraph{Limitations.}
Some reported statistics are difficult to reconcile within any parametric model---for instance, the regression coefficients and bidding frequencies in \citet{kagel1993independent}
cannot both be matched under Gaussian noise, suggesting heteroskedasticity in the original data. Our approach matches key aggregate patterns while acknowledging this limitation.
Comparisons between LLM and human bidding should be interpreted qualitatively. Full parameter values appear in Appendix~\ref{app:mixture}.


\subsection{Unified Metric for all Auctions}\label{sec:method}

\paragraph{Unified deviation metric (applied to all bids).}
Consider an auction \emph{environment} and \emph{format} indexed by $m$.
The index $m$ collects all primitives needed to define the theoretical benchmark, including the number of bidders $N$, the value (or signal) distribution, the information structure (independent private value/common value), and the mechanism rules (e.g., first-price vs.\ second-price, sealed-bid vs.\ clock).
Let $\mathcal I_i$ denote bidder $i$'s information at the moment of bidding under environment $m$ (e.g., $\mathcal I_i=v_i$ in private values, or $\mathcal I_i=s_i$ in common values).
Let $b_{i}$ denote an \emph{observed submitted bid} by bidder $i$ (we apply the definition to \emph{all bids}, not only winning bids), and let $b_m^\star(\mathcal I_i)$ denote the \emph{theoretical benchmark bid} under $m$ given information $\mathcal I_i$ (Bayes--Nash equilibrium in non-strategy-proof formats; truthful/value bidding in strategy-proof formats).

\noindent We define the absolute bid error for a bid with information $\mathcal I_i$ as
\[
d_m(\mathcal I_i)\;\equiv\;\bigl|b_i - b_m^\star(\mathcal I_i)\bigr|.
\]
To compare deviations across formats on a common scale, we normalize by the expected benchmark bid level,
\[
\mu_m^\star \;\equiv\; \mathbb E_m\!\left[b_m^\star(\mathcal I)\right],
\]
where $\mathbb E_m[\cdot]$ denotes expectation with respect to the distribution of the information object $\mathcal I$ induced by environment $m$ (i.e., the value/signal distribution and any randomization in the experimental design).
We then define the \emph{scaled mean absolute deviation} (SMAD) for auction format $m$ as
\[
\Delta_m \;\equiv\; 100 \cdot \frac{\mathbb E_m\!\left[d_m(\mathcal I)\right]}{\mu_m^\star}
\;=\; 100 \cdot \frac{\mathbb E_m\!\left[\left|b-b_m^\star(\mathcal I)\right|\right]}{\mathbb E_m\!\left[b_m^\star(\mathcal I)\right]}.
\]
This metric is interpretable as a percentage of the expected benchmark bid: $\Delta_m=10$ means that, on average, bids deviate from the benchmark by an amount equal to $10\%$ of the expected benchmark bid.


Given a sample of bids $\{(b_k,\mathcal I_k)\}_{k=1}^n$ from format $m$, the plug-in estimator is
\[
\widehat{\Delta}_m
\;=\;
100 \cdot
\frac{\frac{1}{n}\sum_{k=1}^{n}\left|b_k-b_m^\star(\mathcal I_k)\right|}
{\widehat{\mu}_m^\star},
\qquad
\widehat{\mu}_m^\star \equiv \frac{1}{n}\sum_{k=1}^{n} b_m^\star(\mathcal I_k).
\]
When we report a standard error for the (unscaled) mean absolute deviation $\frac{1}{n}\sum_k |b_k-b_m^\star(\mathcal I_k)|$, we propagate it to $\widehat{\Delta}_m$ by scaling with $\widehat{\mu}_m^\star$. We construct $95\%$ confidence intervals using a normal approximation when an analytic standard error is available.


\paragraph{Deterministic benchmarks.}
SMAD requires a single-valued benchmark mapping $\mathcal I \mapsto b_m^\star(\mathcal I)$.
Accordingly, we restrict attention to auction formats where the theoretical prediction is deterministic given the bidder’s information.


\subsection{LLM Data Generation}
\label{sec:data_generation}

We generate synthetic auction data by simulating auctions in which $n$ large language model (LLM) agents bid against one another under fully specified information structures and auction rules. Each experimental \emph{setting} is defined by (i) a \emph{value model} (how bidder values/signals are drawn), (ii) a \emph{mechanism environment} (the action space, allocation and payment rules, and the information revealed during play), and (iii) a \emph{prompt framing} (how the mechanism is described to the bidder). For each setting, we repeat the simulation $K$ times with independent random seeds, producing a dataset of realized values (or signals), bidder actions, allocations, prices, profits, and free-text plans. Here, ``actions'' refers to: a numeric bid in sealed-bid formats, a \texttt{stay}/\texttt{exit} decision at each clock tick in ascending auctions, and an \texttt{increase}/\texttt{hold} decision over a bidder's maximum bid in the eBay proxy-bidding environment.

\paragraph{Single-auction baseline and learning ablation.}
Our main benchmark experiments use \emph{single-auction} simulations: each repetition $k$ is an independent auction instance with freshly drawn values/signals, and bidders do not receive feedback across instances. We adopt this design to minimize reliance on cross-instance learning and long-term adaptation, which remains unreliable for language models \citet{dong2024survey}. Importantly, this choice is distinct from whether a mechanism is \emph{iterative} within an auction (e.g., clock auctions): iterative mechanisms still unfold over multiple within-auction steps (clock ticks), but we do not run repeated auctions with feedback by default. In Appendix~\ref{app:learning}, we run a repeated-auction ablation for sealed-bid auctions in which bidders observe outcomes from prior auction instances; consistent with the motivation above, we find no meaningful learning effects across repetitions.

\paragraph{Value models and numerical calibration.}
We implement three canonical information structures. All parameters (support ranges, bidder count $n$, and bid/price increments) are specified in configuration files and held fixed within a setting. In sealed-bid auctions, bids are restricted to a discrete grid with increment $0.1$; in ascending clock auctions, the clock price increases by increment $0.5$ per tick.

\begin{itemize}
    \item \textbf{Independent Private Values (IPV).} Each bidder's value is the sum of an independent private component and (optionally) a common component. In our baseline IPV calibration, the private component from $[0,49]$, yielding
    \[
          \epsilon_{it} \sim \mathrm{Unif}[0,49], \qquad v_{it}=\epsilon_{it}.
    \]
    We discretize bids with increment $0.1$.

    \item \textbf{Affiliated Private Values (APV).} Values share a common component plus an idiosyncratic private shock. Our APV calibration sets
    \[
        c_t \sim \mathrm{Unif}[0,29], \qquad \epsilon_{it} \sim \mathrm{Unif}[0,20], \qquad v_{it}=c_t+\epsilon_{it},
    \]
    again with bid increment $0.1$. This induces affiliation: higher realized values for one bidder make higher values for others more likely.

    \item \textbf{Common Value (CV).} The item has an underlying common value $V_t$ shared by all bidders, while each bidder observes a noisy private signal. In our CV calibration, the common value is drawn from $[20,29]$ and each bidder observes an idiosyncratically perturbed signal,
    \[
        V_t \sim \mathrm{Unif}[20,29], \qquad \eta_{it} \sim \mathrm{Unif}[-20,20], \qquad s_{it}=V_t+\eta_{it},
    \]
    with bid increment $0.1$. Bidder $i$ conditions on $s_{it}$ when bidding, while realized payoffs depend on $V_t$, generating the winner's-curse selection effect.
\end{itemize}
In our baseline benchmark, $T=1$ so each repetition corresponds to a single auction instance.

\subsubsection{Canonical laboratory mechanisms.}
To mirror standard laboratory benchmarks, we implement \textit{sealed-bid} auctions and \textit{ascending clock} auctions under the value models above.

\paragraph{Sealed-bid auctions.}
In sealed-bid auctions, each bidder submits a single bid $b_{it}$ simultaneously. Let $b_{(1)t}\ge b_{(2)t}\ge b_{(3)t}$ denote the order statistics and let $i^\star_t$ be the highest bidder. The winner pays a price $p_t$ determined by the mechanism:
\begin{itemize}
    \item \textbf{First-price sealed-bid (FP).} $i^\star_t$ wins and pays $p_t=b_{(1)t}$.
    \item \textbf{Second-price sealed-bid (SP).} $i^\star_t$ wins and pays $p_t=b_{(2)t}$.
\end{itemize}
Profits are computed as $\pi_{it}=v_{it}-p_t$ for winners (or $\pi_{it}=V_t-p_t$ in CV) and $0$ for losers.

\paragraph{Ascending clock auctions.}
In an ascending clock auction, the price starts at 0 and increases by a fixed increment $\Delta_{\text{clock}}=0.5$ at each clock tick. At each tick, each bidder chooses \emph{stay} or \emph{exit}. The last remaining bidder wins. We implement the standard second-price logic in the dynamic interface: the winner pays the exit price of the second-to-last bidder.

\paragraph{Framing interventions (mechanism held fixed).}
To isolate the effect of description while holding auction design constant, we implement two alternative framings of the second-price sealed-bid auction:
\begin{itemize}
    \item \textbf{Menu framing.} Bidders are instructed to choose a ``maximum acceptable price.'' They win if the highest rival bid (the market price) is below their chosen maximum, and if they win they pay that market price.
    \item \textbf{Clock framing.} Bidders are told to treat their sealed bid as an automatic exit threshold in a simulated ascending clock: conceptually, the clock rises until bidders ``drop out'' at their thresholds, and the winner pays the second-highest threshold.
\end{itemize}
Both framings implement the same allocation and payment rule as the second-price sealed-bid auction; only the prompt language differs. The precise prompts are listed in Appendix~\ref{app:prompt}.

\subsubsection{Stylized eBay proxy-bidding environment (separate from lab benchmarks).}
Because our eBay-style results section studies a richer, marketplace-inspired microstructure rather than canonical lab mechanisms, we treat eBay simulations as a separate environment. We implement a stylized eBay auction with \emph{proxy bidding} and an optional auction-extension closing rule. Proxy bidding means each bidder maintains a \emph{maximum bid} (their willingness-to-pay cap), and the platform automatically raises the standing bid only as needed to keep the current leader ahead. In each eBay simulation, there are $n=3$ bidders and (unless modified by the closing rule) the auction lasts for $T_{\text{eBay}}=10$ discrete bidding periods. Unless otherwise noted, the eBay environment uses IPV with $v_i \sim \mathrm{Unif}[0,99]$.

\emph{Treatments (reserve price $\times$ closing rule).}
We vary the eBay environment along two design dimensions: whether the auction has a \emph{hidden reserve price} and whether it uses a modified closing rule (auction extension), yielding four treatments:
\begin{itemize}
    \item[\textbf{T1}] Standard eBay auction with proxy bidding.
    \item[\textbf{T2}] eBay auction with a modified (auction extension) closing rule.
    \item[\textbf{T3}] Standard eBay auction with a hidden reserve price.
    \item[\textbf{T4}] eBay with a modified (auction extension) closing rule and a hidden reserve price.
\end{itemize}
A hidden reserve price is a seller-chosen minimum acceptable sale price that is \emph{not disclosed} to bidders; if the highest maximum bid fails to meet the reserve, the item does not sell. For hidden-reserve designs, we consider $r \in \{40, 50, 60\}$.

\emph{Discrete-time implementation of bidding.}
Because LLM interactions are necessarily discrete, we approximate continuous-time bidding using period-based actions. In each period, bidders move in a common-knowledge, randomly permuted order. When it is a bidder's turn, they observe the current auction state (including updates induced by earlier movers in that period) and decide whether to \emph{increase their maximum bid} or \emph{hold}. In the final period, bidding is simultaneous: bidders submit their final decisions without observing others' last-period actions, capturing the idea that, in continuous time, a bidder cannot know whether another bid will arrive just before closing.

\emph{Modified closing rule (auction extension).}
Under the modified closing rule, if a higher maximum bid is placed in the final period such that a new bidder becomes the winner, the auction is extended by one additional period. This process repeats until a final period occurs in which no further bids are made.

\emph{Abstraction and prompting.}
This stylized eBay framework abstracts away from additional platform features such as shipping fees and Buy-It-Now options. Moreover, unlike the explicit framing interventions in Section~\ref{sec:Intervention}, we keep the eBay prompts deliberately \emph{minimal} and descriptive: the goal is to study how LLMs behave under richer market microstructure without steering them toward any particular strategy through didactic instruction. The complete simulation process is described in Appendix~\ref{app:ebay}, and the corresponding prompts are listed in Appendix~\ref{app:prompt}.

\subsubsection{Simulation protocol.}
Each experiment proceeds as follows, repeated for $K$ independent repetitions:
\begin{enumerate}
    \item \textbf{Value realization.} For each repetition $k$, we draw values (IPV/APV) or signals (CV) for all bidders under a fixed random seed. In our baseline benchmark experiments, each repetition corresponds to a single auction instance.
    \item \textbf{Prompted action.} We construct an individualized prompt for each bidder containing (i) the mechanism rules (and framing, if applicable), (ii) the bidder's current value (or signal), and (iii) when applicable, a feedback transcript summarizing prior auction instances (only in the learning ablation). We request a structured response with explicit tags:
\begin{quote}\small
\verb|<PLAN> ... </PLAN>| \\
\verb|<ACTION> ... </ACTION>|
\end{quote}
where \verb|<ACTION>| contains a numeric bid (sealed-bid), a stay/exit decision (clock), or an increase/hold decision over the bidder's maximum bid (eBay).
    \item \textbf{Parallel elicitation and parsing.} We query the LLM for each bidder using separate API calls in parallel, and parse the action from the \verb|<ACTION>| tag using a deterministic regex-based extractor.
    \item \textbf{Outcome computation.} Given the submitted actions, we determine the winner and payments according to the mechanism rules and compute bidder profits.
\end{enumerate}

Unless otherwise noted, our baseline settings use $n=3$ bidders and we repeat each configuration $K$ times with independent seeds to ensure stable distributional estimates.


\subsection{Model Selection}

To run the simulation, we evaluate the following five LLMs under identical prompts:  we choose 3 most popular foundation models from different providers: OpenAI’s gpt-4o (2024-08-06),  Anthropic's claude-3-5-haiku (2024-10-22) and Google's gemini-2.0-flash (2025-0205), one reasoner from the canonical LLM: gpt-5-mini (2025-08-07) ; and one widelyused open-sourced model: Meta-Llama-3-8B-Instruct.
All models are prompted with the survey instrument described in the previous section.
For non-reasoning models (e.g., gpt-4o, Llama), we set the temperature to $temp = 0.5$,  which is the most common setting in LLM simulation\citep{horton2023large, Horton2024EDSL}. For reasoning models (GPT-5-mini), temperature is ignored by the API.
For our baseline gpt-4o, we also ablate set the temperature to  $ temp\in \{ 0.1, 1.0\}$.








